\documentclass[11pt,a4paper]{article}

\usepackage[a4paper,margin=2cm]{geometry}
\usepackage{fontspec}
\usepackage{polyglossia}
\usepackage{array}
\usepackage{longtable}
\usepackage{tabularx}
\usepackage{enumitem}
\usepackage{xcolor}
\usepackage{hyperref}
\usepackage{titlesec}
\usepackage{fancyhdr}
\usepackage{setspace}

\setmainlanguage{greek}
\setotherlanguage{english}
\setmainfont{Arial}
\newfontfamily\greekfont{Arial}

\definecolor{ink}{HTML}{163447}
\definecolor{soft}{HTML}{5F7888}
\definecolor{line}{HTML}{D6E1E8}
\definecolor{accent}{HTML}{1F6F7E}
\definecolor{warnbg}{HTML}{FFF7E4}
\definecolor{warnline}{HTML}{F3B533}

\hypersetup{
  colorlinks=true,
  linkcolor=accent,
  urlcolor=accent,
  pdftitle={Οδηγίες Χρήσης Εφαρμογής Field Operations},
  pdfauthor={Codex}
}

\titleformat{\section}{\Large\bfseries\color{ink}}{\thesection.}{0.6em}{}
\titleformat{\subsection}{\large\bfseries\color{ink}}{\thesubsection.}{0.6em}{}

\setlength{\parindent}{0pt}
\setlength{\parskip}{0.55em}
\setlist[itemize]{itemsep=0.35em, topsep=0.35em}
\setlist[enumerate]{itemsep=0.35em, topsep=0.35em}
\renewcommand{\arraystretch}{1.35}

\pagestyle{fancy}
\fancyhf{}
\fancyhead[L]{\textcolor{soft}{Field Operations Manual}}
\fancyhead[R]{\textcolor{soft}{\thepage}}
\renewcommand{\headrulewidth}{0.3pt}
\renewcommand{\headrule}{\hbox to\headwidth{\color{line}\leaders\hrule height \headrulewidth\hfill}}

\newcommand{\infobox}[1]{
  \vspace{0.5em}
  \fcolorbox{line}{white}{
    \parbox{\dimexpr\linewidth-2\fboxsep-2\fboxrule\relax}{
      #1
    }
  }
  \vspace{0.5em}
}

\newcommand{\warningbox}[1]{
  \vspace{0.5em}
  {\setlength{\fboxrule}{1pt}
  \fcolorbox{warnline}{warnbg}{
    \parbox{\dimexpr\linewidth-2\fboxsep-2\fboxrule\relax}{
      \textbf{Σημαντικό:} #1
    }
  }}
  \vspace{0.5em}
}

\begin{document}

{\small\textcolor{soft}{Field Operations Manual}}

{\huge\bfseries\color{ink} Οδηγίες Χρήσης για Υπάλληλο}

Το παρόν εγχειρίδιο εξηγεί βήμα βήμα πώς χρησιμοποιείται η εφαρμογή διαχείρισης εργασιών πεδίου. Είναι γραμμένο για καθημερινή χρήση από υπάλληλο ή συνεργάτη που πρέπει να βλέπει εργασίες, να τις αναλαμβάνει, να τις ενημερώνει και να τις προωθεί σωστά μέσα στο workflow.

\vspace{0.7em}

\begin{tabularx}{\linewidth}{|>{\raggedright\arraybackslash}p{0.29\linewidth}|>{\raggedright\arraybackslash}p{0.33\linewidth}|>{\raggedright\arraybackslash}X|}
\hline
\textbf{Ρόλοι συστήματος} & \textbf{Βασική λειτουργία} & \textbf{Workflow} \\
\hline
Admin 1, Συνεργάτης 1, Συνεργάτης 2 & Διαχείριση και παρακολούθηση τεχνικών εργασιών πεδίου & Ανατέθηκε $\rightarrow$ Προγραμματισμένη $\rightarrow$ Σε εξέλιξη $\rightarrow$ Για επικύρωση $\rightarrow$ Ολοκληρωμένη \\
\hline
\end{tabularx}

\warningbox{Η εφαρμογή είναι prototype χωρίς database. Τα δεδομένα αποθηκεύονται τοπικά στον browser του χρήστη. Αυτό σημαίνει ότι χρησιμοποιείται για δοκιμή, παρουσίαση, έλεγχο ροής και όχι για τελική παραγωγική χρήση.}

\section{Τι κάνει η εφαρμογή}

Η εφαρμογή χρησιμοποιείται για να οργανώνει τεχνικές εργασίες πεδίου. Κάθε εργασία έχει στοιχεία έργου, διεύθυνση, υπεύθυνο συνεργάτη, προγραμματισμό, φωτογραφίες, αρχεία, υλικά, ιστορικό και τεχνικά metadata.

Ο βασικός στόχος είναι να μην υπάρχει μόνο μια απλή λίστα εργασιών, αλλά πλήρης έλεγχος του κύκλου ζωής κάθε εργασίας από τη δημιουργία μέχρι την έγκριση και το κλείσιμο.

\section{Ποιοι ρόλοι υπάρχουν}

\begin{longtable}{|>{\raggedright\arraybackslash}p{0.2\linewidth}|>{\raggedright\arraybackslash}p{0.72\linewidth}|}
\hline
\textbf{Ρόλος} & \textbf{Τι μπορεί να κάνει} \\
\hline
Admin 1 & Βλέπει όλες τις εργασίες, δημιουργεί νέες, ορίζει ποιοι συνεργάτες έχουν πρόσβαση, κάνει αλλαγές, εγκρίνει ή απορρίπτει. \\
\hline
Συνεργάτης 1 & Βλέπει μόνο τις εργασίες που του επιτρέπει ο admin ή όσες έχει ήδη αναλάβει ο ίδιος. \\
\hline
Συνεργάτης 2 & Έχει την ίδια λογική με τον Συνεργάτη 1 και δουλεύει μόνο μέσα στο δικό του επιτρεπόμενο scope. \\
\hline
\end{longtable}

\infobox{Η επιλογή ρόλου και χρήστη γίνεται από το πάνω δεξιά μέρος της εφαρμογής. Δεν υπάρχει login με κωδικό στο prototype.}

\section{Κεντρικές ενότητες της εφαρμογής}

\begin{itemize}
  \item \textbf{Dashboard}: Δείχνει πόσες εργασίες υπάρχουν σε κάθε κατάσταση.
  \item \textbf{Εργασίες}: Δείχνει τον πίνακα με όλες τις ορατές εργασίες και επιτρέπει αναζήτηση και φιλτράρισμα.
  \item \textbf{Task Detail}: Είναι η πλήρης καρτέλα της εργασίας με tabs, workflow actions και ιστορικό.
  \item \textbf{Export PDF}: Εξάγει αναφορά με όλες τις ανοιχτές εργασίες που μπορεί να δει ο τρέχων χρήστης.
\end{itemize}

\section{Πώς ξεκινά ο χρήστης}

\begin{enumerate}
  \item Από το επάνω menu επιλέγει αν θα δουλέψει ως \textbf{Admin} ή ως \textbf{Partner}.
  \item Από το δεύτερο dropdown επιλέγει συγκεκριμένο χρήστη, π.χ. \textbf{Admin 1} ή \textbf{Συνεργάτης 1}.
  \item Στο dashboard βλέπει αμέσως πόσες εργασίες υπάρχουν ανά στάδιο για τον τρέχοντα χρήστη.
\end{enumerate}

\section{Πώς δημιουργεί εργασία ο Admin}

\begin{enumerate}
  \item Ο admin πατά το κουμπί \textbf{Νέα εργασία}.
  \item Ανοίγει φόρμα δημιουργίας με τα βασικά πεδία.
  \item Συμπληρώνει τίτλο, είδος, project, project ID, SR / BID, team, διεύθυνση, πόλη και σημειώσεις.
  \item Ορίζει προγραμματισμό αν υπάρχει ήδη γνωστή ημερομηνία και ώρα.
  \item Επιλέγει ποιοι συνεργάτες επιτρέπεται να δουν την εργασία.
  \item Πατά \textbf{Δημιουργία}.
\end{enumerate}

Μετά τη δημιουργία, το σύστημα παράγει αυτόματα:

\begin{itemize}
  \item Task ID
  \item Created at / Created by
  \item Updated at / Updated by
  \item Αρχική εγγραφή στο History
  \item Κενά πεδία για φωτογραφίες, αρχεία και υλικά
\end{itemize}

\section{Τι σημαίνει κάθε βασικό πεδίο}

\begin{longtable}{|>{\raggedright\arraybackslash}p{0.28\linewidth}|>{\raggedright\arraybackslash}p{0.64\linewidth}|}
\hline
\textbf{Πεδίο} & \textbf{Επεξήγηση} \\
\hline
Τίτλος & Το βασικό όνομα της εργασίας. \\
\hline
Είδος & Αυτοψία, εγκατάσταση ή επισκευή. \\
\hline
Project & Το έργο στο οποίο ανήκει η εργασία. \\
\hline
Project ID & Ο κωδικός αναφοράς του έργου. \\
\hline
SR / BID & Εσωτερικός ή εξωτερικός κωδικός ticket / αιτήματος. \\
\hline
Πόρος / Team & Η ομάδα ή το συνεργείο που σχετίζεται με την εργασία. \\
\hline
Επιτρεπόμενοι partners & Ποιοι συνεργάτες θα μπορούν να δουν και να αναλάβουν την εργασία. \\
\hline
Αναλήφθηκε από partner & Ποιος συνεργάτης τελικά την έχει αναλάβει. \\
\hline
Ημ/νία έναρξης & Πότε έχει προγραμματιστεί να ξεκινήσει. \\
\hline
Ημ/νία λήξης & Πότε ολοκληρώθηκε ή αναμένεται να ολοκληρωθεί. \\
\hline
Σημειώσεις & Ελεύθερο πεδίο με οδηγίες, παρατηρήσεις ή ενημερώσεις. \\
\hline
\end{longtable}

\section{Πώς βλέπει και αναλαμβάνει εργασία ο συνεργάτης}

\begin{enumerate}
  \item Ο χρήστης επιλέγει ρόλο \textbf{Partner}.
  \item Επιλέγει τον εαυτό του, π.χ. \textbf{Συνεργάτης 1}.
  \item Πηγαίνει στη σελίδα \textbf{Εργασίες}.
  \item Βλέπει μόνο τις εργασίες που επιτρέπεται να δει.
  \item Πατά πάνω στη γραμμή της εργασίας για να μπει στην καρτέλα της.
  \item Αν η εργασία είναι διαθέσιμη, πατά \textbf{Ανάληψη εργασίας}.
\end{enumerate}

\warningbox{Αν ο συνεργάτης δεν είναι μέσα στους επιτρεπόμενους χρήστες, δεν θα βλέπει την εργασία καθόλου.}

\section{Πώς προχωρά το workflow}

\begin{longtable}{|>{\raggedright\arraybackslash}p{0.2\linewidth}|>{\raggedright\arraybackslash}p{0.38\linewidth}|>{\raggedright\arraybackslash}p{0.32\linewidth}|}
\hline
\textbf{Κατάσταση} & \textbf{Τι σημαίνει} & \textbf{Ποιος τη μετακινεί} \\
\hline
Ανατέθηκε & Η εργασία υπάρχει αλλά δεν έχει ολοκληρωθεί η ενεργή εκτέλεση. & Admin δημιουργεί την εργασία. \\
\hline
Προγραμματισμένη & Υπάρχει συνεργάτης και χρονικό παράθυρο. & Admin ή συνεργάτης μετά την ανάληψη. \\
\hline
Σε εξέλιξη & Η εργασία εκτελείται στο πεδίο. & Ο συνεργάτης πατά \textbf{Έναρξη εργασίας}. \\
\hline
Για επικύρωση & Η εκτέλεση ολοκληρώθηκε και αναμένει έλεγχο από admin. & Ο συνεργάτης πατά \textbf{Αποστολή για επικύρωση}. \\
\hline
Ολοκληρωμένη & Η εργασία εγκρίθηκε και έκλεισε. & Ο admin πατά \textbf{Έγκριση ολοκλήρωσης}. \\
\hline
\end{longtable}

Αν ο admin διαπιστώσει πρόβλημα, μπορεί να πατήσει \textbf{Απόρριψη και επιστροφή}. Τότε η εργασία επιστρέφει σε κατάσταση \textbf{Σε εξέλιξη}.

\section{Τι κάνει ο χρήστης μέσα στην καρτέλα εργασίας}

\begin{longtable}{|>{\raggedright\arraybackslash}p{0.22\linewidth}|>{\raggedright\arraybackslash}p{0.7\linewidth}|}
\hline
\textbf{Tab} & \textbf{Χρήση} \\
\hline
Κύριος & Τα βασικά στοιχεία της εργασίας, το planning, οι συνεργάτες και οι σημειώσεις. \\
\hline
Φωτογραφίες & Ανέβασμα φωτογραφιών σε κατηγορίες πριν, μετά, εξοπλισμός, καλωδίωση. \\
\hline
Αρχεία & Ανέβασμα PDF και λοιπών υποστηρικτικών αρχείων. \\
\hline
Υλικά & Καταχώρηση υλικών που χρησιμοποιήθηκαν. \\
\hline
Floors & Πληροφορίες για ορόφους, μονάδες και risers του κτιρίου. \\
\hline
Health \& Safety & Checklist ασφάλειας, παρατηρήσεις και σημεία προσοχής. \\
\hline
Ιστορία & Δείχνει ποιος έκανε τι και πότε. \\
\hline
Σύστημα & Τεχνικά metadata όπως task ID, timestamps και internal flags. \\
\hline
\end{longtable}

\section{Πώς ανεβάζει ο συνεργάτης τεκμηρίωση}

\begin{enumerate}
  \item Μπαίνει στο tab \textbf{Φωτογραφίες} και επιλέγει την κατάλληλη κατηγορία.
  \item Ανεβάζει όσες φωτογραφίες χρειάζονται.
  \item Μπαίνει στο tab \textbf{Αρχεία} και ανεβάζει report ή PDF.
  \item Αν χρησιμοποιήθηκαν υλικά, τα περνά στο tab \textbf{Υλικά}.
  \item Συμπληρώνει τυχόν παρατηρήσεις στο \textbf{Health \& Safety}.
\end{enumerate}

\section{Τι πρέπει να κάνει ο Admin πριν εγκρίνει}

\begin{itemize}
  \item Να ελέγξει αν η εργασία είναι όντως στη φάση \textbf{Για επικύρωση}.
  \item Να επιβεβαιώσει ότι υπάρχουν σημειώσεις και σωστό planning.
  \item Να ελέγξει φωτογραφίες και αρχεία.
  \item Να ελέγξει ότι η καρτέλα δεν έχει εμφανή ελλείμματα.
  \item Να διαβάσει το ιστορικό για να δει τι αλλαγές έγιναν.
\end{itemize}

Αν όλα είναι σωστά, ο admin πατά \textbf{Έγκριση ολοκλήρωσης}. Αν όχι, γράφει σχόλιο και πατά \textbf{Απόρριψη και επιστροφή}.

\section{Πώς γίνεται το Export PDF}

\begin{enumerate}
  \item Από το αριστερό menu επιλέγουμε \textbf{Export PDF}.
  \item Ανοίγει αναφορά με όλες τις ανοιχτές εργασίες που μπορεί να δει ο τρέχων χρήστης.
  \item Πατάμε \textbf{Εκτύπωση / Save as PDF}.
  \item Στο παράθυρο εκτύπωσης του browser επιλέγουμε αποθήκευση ως PDF.
\end{enumerate}

\infobox{Ο admin εξάγει όλες τις ανοιχτές εργασίες. Ο συνεργάτης εξάγει μόνο όσες επιτρέπεται να βλέπει.}

\section{Τι κάνει το Reset demo}

Το κουμπί \textbf{Reset demo} επαναφέρει τα mock δεδομένα της εφαρμογής στην αρχική τους κατάσταση. Χρησιμοποιείται όταν θέλουμε να ξαναδείξουμε την εφαρμογή από την αρχή ή όταν θέλουμε να καθαρίσουμε τις δοκιμές.

\section{Συχνά λάθη που πρέπει να αποφεύγονται}

\begin{itemize}
  \item Να ξεχνά ο admin να δώσει πρόσβαση στον σωστό συνεργάτη.
  \item Να προχωρά εργασία χωρίς να έχουν ανέβει βασικά αποδεικτικά στοιχεία.
  \item Να γίνεται έγκριση χωρίς έλεγχο των tabs φωτογραφιών, αρχείων και ιστορικού.
  \item Να γίνεται Reset demo χωρίς να γνωρίζει ο χρήστης ότι θα χαθούν οι τοπικές δοκιμαστικές αλλαγές.
\end{itemize}

\section{Σύντομη καθημερινή ροή χρήσης}

\begin{enumerate}
  \item Ο admin δημιουργεί εργασία.
  \item Ο admin επιλέγει ποιοι συνεργάτες μπορούν να τη δουν.
  \item Ο συνεργάτης μπαίνει και την αναλαμβάνει.
  \item Ο συνεργάτης ξεκινά την εργασία.
  \item Ο συνεργάτης ενημερώνει στοιχεία, φωτογραφίες, αρχεία και υλικά.
  \item Ο συνεργάτης τη στέλνει για επικύρωση.
  \item Ο admin την εγκρίνει ή την επιστρέφει για διόρθωση.
\end{enumerate}

\vspace{1em}
\hrule
\vspace{0.6em}

{\small\textcolor{soft}{Έγγραφο χρήσης για το prototype “Field Operations Control”. Το manual αυτό αντιστοιχεί στην τρέχουσα frontend έκδοση της εφαρμογής χωρίς backend και χωρίς κοινή βάση δεδομένων.}}

\end{document}
